\documentclass{article}
\usepackage[T1]{fontenc}
\usepackage[utf8]{inputenc}
\usepackage[polish]{babel}
\usepackage{caption}
\usepackage{enumitem}
\usepackage{amsmath}
\usepackage{graphicx}
\usepackage{url}


\title{Mercury}
\author{Karol Krawczykiewicz, Grzegorz Rogoziński, Jan Król, Piotr Maszczak}
\date{Październik 2023 - Styczeń 2024}

\begin{document}

\begin{figure}
    \centering
    \includegraphics[width=0.2\textwidth]{pics/logo.png}
\end{figure}

\maketitle
\section{Informacje ogólne}
\textbf{Nazwa projektu:} Mercury
\\
\\
\textbf{Technologie:} Express, MongoDB, Neo4j, React, Redux, Sockets, TailwindCSS, TypeScript, WebRTC
\\
\\
\textbf{Protokoły:} HTTP, UDP, ICE, SDP
\\
\\
\textbf{Repozytorium GitHub:}
\url{https://github.com/Karol-2/Mercury-Project}

\section{Opis projektu}
Mercury jest aplikacją webową, która zapewnia komunikację tekstową oraz na żywo. Aplikacja jest dostosowana do przeglądarek internetowych na komputerach i telefonach, umożliwiając użytkownikom współpracę na różnych urządzeniach. Ponadto, system zapewnia możliwości wyszukiwania, dodawania i usuwania znajomych.


\newpage
\section{Opis działania}
\subsection{Połączenia}

Aplikacja React łączy się z API backendu. Ten zaś, posiada zabezpieczone połączenie z 2 bazami danych. System składa się z 4 elementów:
\begin{enumerate}
  \item Frontend - działa na \url{http://localhost:5173}
  \item Backend - działa na \url{http://localhost:5000}
  \item Baza Neo4j - przechowuje dane użytkowników,
  działa na \url{bolt://localhost:7687} do komunikacji z bazą danych
  i \url{http://localhost:7474} do testowania kwerend i wizualizacji danych
  \item Baza MongoDB - przechowuje dane o chatach,  działa na porcie 27017
\end{enumerate}
API Backendu udostępnia ścieżki dotyczące autoryzacji, obsługi czatów, edycji relacji i obsługi użytkowników. Zapytania są zabezpieczone tokenami.

\subsection{Uruchomienie}
Aplikację można uruchomić na dwa sposoby:
\begin{enumerate}
    \item Uruchamiając backend i frontend lokalnie, używając narzędzia npm, i łącząc z lokalnymi wersjami baz danych.
    \item Używając kontenerów Docker, dzięki plikowi \textbf{compose.yml} z konfiguracją.
\end{enumerate}

\newpage
\section{Modele baz danych}
\subsection{Model użytkownika}
\begin{figure}[h]
    \centering
    \includegraphics[width=0.5\textwidth]{pics/user_model.png}
\end{figure}

Model użytkownika składa się z ośmiu pól: \nolinebreak
\begin{enumerate}[itemsep=0mm]
  \item \textbf{id} - identyfikator
  \item \textbf{first\_name} - imię
  \item \textbf{last\_name} - nazwisko
  \item \textbf{country} - dwuliterowy kod państwa
  \item \textbf{profile\_picture} - obraz zaszyfrowany base64
  \item \textbf{mail} - adres email
  \item \textbf{password} - zaszyfrowane hasło
\end{enumerate}

\newpage
\subsection{Relacje między użytkownikami}
\begin{figure}[h]
    \centering
    \includegraphics[width=0.8\textwidth]{pics/friends_nodes.png}
    \caption*{Użytkownicy z dwustronną relacją}
\end{figure}

Jeżeli dwaj użytkownicy są przyjaciółmi, istnieje miedzy nimi relacja \\
\textbf{IS\_FRIENDS\_WITH}. Obowiązuje ona w dwie strony.

\begin{figure}[h]
    \centering
    \includegraphics[width=0.8\textwidth]{pics/send_request.png}
    \caption*{Użytkownicy z jednostronną relacją}
\end{figure}

W sytuacji gdy jeden użytkownik wyśle zaproszenie do drugiego, nawiązuje się między nimi relacja \textbf{SENT\_INVITE\_TO}. Osoba, od której wychodzi strzałka, wysłała prośbę o dodanie do osoby, przy której jest grot strzałki.
Gdy zostaje ono zaakceptowane, usuwane są dotychczasowe relacje i nadawana jest dwustronna relacja \textbf{IS\_FRIENDS\_WITH}.

\begin{figure}[ht]
    \centering
    \includegraphics[width=0.8\textwidth]{pics/no_friends_nodes.png}
    \caption*{Użytkownicy bez relacji}
\end{figure}

W przypadku anulowania prośby lub usunięcia ze znajomych, wszystkie relacje między dwoma użytkownikami zostają usunięte.

\subsection{Model chatu}
\begin{figure}[h]
    \centering
    \includegraphics{pics/model_mongo.png}
\end{figure}

Model chatu zawiera dane o wiadomościach tekstowych przesyłanych między użytkownikami. Składają się na niego pola: \textbf{id} (id wiadomości), \textbf{authorId} (id nadawcy), \textbf{receiverId} (id odbiorcy), \textbf{content} (wiadomość), \textbf{created\_date} (data wysłania). 

\newpage
\section{Najważniejsze systemy}
\subsection{System rejestracji i logowania}

Rejestracja nowego użytkownika odbywa się po wybraniu opcji Register na ekranie głównym, lub poprzez wejście na odpowiedni endpoint: \\
\url{http://localhost:5173/register}

\begin{figure}[h]
    \centering
    \includegraphics[height=0.8\textwidth]{pics/register.png}
    \caption*{Formularz rejestracyjny}
\end{figure}

Użytkownik, aby móc się zarejestrować, musi podać swoje dane, które są w odpowiedni sposób walidowane. Reguły w formularzu rejestracji:
\begin{enumerate}
    \item Imię i nazwisko z przynajmniej 2 znakami
    \item Kod państwa o dokładnie 2 literach
    \item Brak zajętego konta o tym samym mailu
    \item Email o poprawnym formacie
    \item Hasło z przynamniej 8 znakami
\end{enumerate}

Zdjęcie profilowe nie jest wymagane. W przypadku braku wprowadzenia własnego zdjęcia, system sam wprowadza domyślne zdjęcie. Wybór zdjęcia jest widoczny jako miniaturka, również z możliwością usunięcia.

Dane z formularza są walidowane w czasie rzeczywistym przy użyciu biblioteki Zod
i React hooka useForm. Po wypełnieniu formularza i wciśnięciu przycisku 'Register'
przesłany obrazek (lub w przypadku braku, to obrazek domyślny) jest przekodowywany
na format base64. Hasło jest szyfrowane za pomocą biblioteki bcrypt. Następnie
sprawdzane jest, czy podany maila z formularza nie jest zajęty w systemie.
Jeśli jest zajęty, to wyświetlony jest komunikat z tym związany, formularz
nie jest czyszczony w celu poprawienia maila. W przypadku poprawności danych,
system zapisuje nowego użytkownika oraz następuje zmiana strony
na logowanie na endpointcie: \url{http://localhost:5173/login}

\begin{figure}[h]
    \centering
    \includegraphics[height=0.8\textwidth]{pics/login.png}
    \caption*{Formularz logowania}
\end{figure}

Formularz logowania wymaga podania maila oraz hasła. Po poprawnym podaniu danych
jest tworzona sesja z parą tokenów:
\begin{enumerate}
  \item access token - token o krótkim czasie ważności (900 sekund - 15 minut)
  przechowywany w sessionStorage
  \item refresh token - token o długim czasie ważności (1209600 sekund - 2 tygodnie)
  przechowywany na serwerze w ramach ciasteczka secure, httpOnly (przechowywany
  po stronie serwera)
\end{enumerate}
Access token pozwala na dostęp do zasobów. Każdy, kto go posiada ma dobrowolny
dostęp do zasobów danego użytkownika. Z tego powodu ważny jest jego krótki czas ważności.
W przypadku utraty ważności klient używa refresh token'a do otrzymania nowego
access token'a. Ta operacja jest bezpieczna, ponieważ refresh token
jest przechowywany od strony serwera.

Od strony backendu są używane biblioteki cookie-parser i jsonwebtoken
a od strony frontendu biblioteka js-cookie.
Po zalogowaniu się następuje przekierowanie na URL \url{http://localhost:5173/profile}
do którego ma się tylko dostęp podczas sesji. Zawiera on wszystkie dane, które zostały
wcześniej podane w rejestracji.

\newpage
\subsection{System szukania znajomych}
Wyszukiwanie znajomych odbywa się na endpointcie \url{http://localhost:5173/search}.
Po wpisaniu w pasek wyszukiwania danej frazy i zatwierdzeniu jest pokazywane
10 pierwszych wyników, które są najbardziej zbliżone do tej frazy
na podstawie imienia i nazwiska. W celu wyszukiwania podobieństwa jest wykorzystywany
własny algorytm oparty o algorytmie odległości Levenshteina.
Każda osoba na liście zawiera przycisk informujący relację z danym kontaktem.

\begin{figure}[h]
    \centering
    \includegraphics[height=0.6\textwidth]{pics/search.png}
    \caption*{Okno wyszukiwania znajomych}
\end{figure}

Od strony backendu jest wykorzystany vector search index.
Imię i nazwisko są przechowywane jako wektor wartości.
Jeżeli jedna para imienia i nazwiska jest podobna do drugiej
to odległość między dwoma wektorami powinna być mała. W tym przypadku metryka
wykorzystywana do porównania wektorów to cosine similarity.

Do zamiany tekstu na wektor (word embedding) jest wykorzystywany własny algorytm.
Miał on do zrealizowania cele:
\begin{enumerate}[itemsep=0mm]
  \item nie powinien wykorzystywać sztucznej inteligencji (brak danych testowych
  i długiego trenowania),
  \item słowa z literami, które są blisko siebie na układzie QWERTY
  powinny być podobne,
  \item słowa z literami zamienionymi miejscami powinny być podobne,
  \item słowa z brakującymi lub dodanymi literami powinny być podobne,
  \item część słowa powinna być podobna do całego słowa.
\end{enumerate}
Pierwszym krokiem algorytmu jest zamiana każdej litery na wartość z zakresu
$[-1, 1]$. Wartość wskazuje, po której stronie klawiatury znajduje się klawisz,
na przykład litera Q ma wartość -1, P ma wartość 1.
Porównanie dwóch słów o tej samej długości, ale np. ze zmienioną literą
A na Q będzie dawało mały wynik a na przykład z A na L - większy.

Drugim krokiem jest użycie interpolacji liniowej do rozszerzenia wektorów
do tej samej, stałej długości. Słowo dłuższe albo krótsze o jedną literę
nadal będą się dobrze pokrywać i dadzą bliski wynik.

Trzecim krokiem jest połączenie poprzedniego wyliczonego wektora
z wektorem z posortowanymi wartościami. Przy porównywaniu dwóch słów,
w tym jednym z zamienionymi literami pierwszy wektor da niższy wynik
ale drugi wyższy, co zwiększa odporność algorytmu na zamianę liter.

\newpage
\subsection{System dodawania do znajomych}
Po wybraniu kontaktu, który ma zostać dodany do sieci znajomych, wysyłany jest Request i tworzona relacja \textbf{SENT\_INVITE\_TO} w bazie Neo4j. Przychodzące zaproszenia i znajomych można sprawdzić w zakładce \textit{Friends} lub na endpointcie \url{http://localhost:5173/friends}

\begin{figure}[h]
    \centering
    \includegraphics[width=1\textwidth]{pics/friends.png}
    \caption*{Ekran listy znajomych i przychodzących zaproszeń}
\end{figure}

Zaproszenia do znajomych są wczytywane cały czas na żywo, poprzez wysyłanie requesta do bazy danych o sprawdzeniu relacjach z zalogowanym użytkownikiem. 
% Endpointy wysyłane na bazę danych dla przykładowego id:
% \begin{enumerate}
%     \item /users/userId/friend-requests - (GET) wczytanie wszystkich zaproszeń do znajomych
%     \item /users/userId/friends - (GET) wczytanie wszystkich obecnych znajomych
%     \item /users/userId/remove/friendId - (DELETE) odrzucenie zaproszenia, usunięcie wszelkich relacji pomiędzy użytkownikiem o id 'userId' oraz 'friendId'
%     \item /users/userId/accept/friendId - (POST) - zaakceptowanie zaproszenia, stworzenie relacji \textbf{IS\_FRIENDS\_WITH} z użytkownikiem o id 'userId' z 'friendId'
% \end{enumerate}

\subsection{System wideo rozmowy}
Użycie WebRTC \textit{(Web Real-Time Communication)} pozwala nawiązać połączenie typu \textit{peer to peer} między dwoma użytkownikami i wymieniać się na bieżąco sygnałem audio i wideo. To rozwiązanie nie potrzebuje serwera, co zmniejsza opóźnienie między użytkownikami. WebRTC transportuje dane z użyciem protokołów UDP, który jest ceniony za swoją szybkość
\\\\
 Schemat komunikacji
 \begin{enumerate}
     \item Gdy Użytkownik1 chce połączyć się z Użytkownikiem2, wysyła on odpowiednią wiadomość o chęci dołączenia.
     \item  Użytkownik2 chcąc połączyć się, akceptuje 'ofertę' i wysyła Użytkownikowi1 swoje informacje.
     \item  Po tym, gdy użytkownicy wymienią się danymi, nawiązuje się połączenie, każdy użytkownik zna SDP drugiego.
     \item Używając metody ICE, każdy użytkownik uderza do \textit{stun server} by uzyskać swój publiczny adres IP.
     \item Gdy \textit{stun server} odpowie, odebrane dane są poprzez użytkownika transportowane do drugiego. Tak samo działa to u drugiego użytkownika.
     \item Kiedy dane znajdą wspólną drogę komunikacji, połączenie jest już gotowe i mogą być przesyłane informacje w obie strony.
 \end{enumerate}

Protokół WebRTC zajmuje się tylko przesyłaniem strumieni audio i wideo.
Do przesyłania komunikatów o dołączeniu wykorzystuje się inny protokół - w tym
projekcie komunikacja odbywa się przez web sockety.

Organizacją procesu zajmuje się ICE\textit{(Interactive Connectivity Establishment)} sygnaling.

 W danych, które są zawierane podczas dołączenia do rozmowy między użytkownikami, znajduje się protokół SDP \textit{(Session Decription Protocol)}, który zawiera informacje typu: kodek, adres, typ nośnika, dane audio, dane wideo. Użytkownicy również wymieniają się '\textit{ICE candidates}', którym jest publiczny adres IP i port, który przyjmuje dostarczane dane.

\subsection{System chatu}

Chat ze znajomymi jest dostępny poprzez wybranie zielonego przycisku chatu obok znajomego w liście dostępnych znajomych. Nowy chat jest tworzony wraz z przypisanym id. Aby wysłać wiadomość, należy wcisnąć Enter.

Wiadomości wysyłane są w czasie rzeczywistym, a cała konwersacja pomiędzy użytkownikami jest zapisywana do bazy MongoDB. Do komunikacji wykorzystwane są gniazda WebSocket. Historia wiadomości jest dostępna pod endpointem: \textbf{/chat/userId/friendId}

\begin{figure}[h]
    \centering
    \includegraphics[width=0.8\textwidth]{pics/messages.png}
    \caption*{Przykładowy ekran chatu}
\end{figure}


\section{Źródła}
\begin{itemize}
    \item Dokumentacja Framer-motion: \url{https://www.framer.com/motion/}
    \item Dokumentacja Zod: \url{https://zod.dev/}
    \item Dokumentacja Bcrypt: \url{https://www.npmjs.com/package/bcrypt}
    \item Dokumentacja Express: \url{https://expressjs.com/en/api.html}
    \item Artykuł Praca Zespołowa na Github: \\
    \url{https://www.freecodecamp.org/news/how-to-use-git-and-github-in-a-team-like-a-pro/}
    \item Kurs WebRTC React: \\
    \url{https://www.udemy.com/course/mastering-webrtc-part-2-real-time-video-and-screen-share/}
    \item Kurs WebRTC: \url{https://www.youtube.com/watch?v=QsH8FL0952k}
\end{itemize}

\end{document}
